% Patch record for C-69 -- resolve the Roll RMS contradiction Sam flagged in PR #17 (C-65).
% Applied 2026-10-31. Author decision: scope the claim and explain the table, exclude no data.
%
% THE CONTRADICTION AS REPORTED
% results.tex (transition-stability paragraph, written by C-61 S-05) said the mean Roll RMS "stays at
% or below 1.59 deg ... across these suites and both variable-topography terrains", while
% tab:fsm_comparison reports Inclined Slope / Neural Roll RMS = 6.751 +/- 7.160 deg (N=3). A reviewer
% comparing the two reads a contradiction.
%
% WHAT THE DATA SHOWS (investigated 2026-10-31, from the archived per-trial logs)
% There are TWO independent batteries of the same two terrain scenarios, and the manuscript was citing
% both as if they were one. Roll RMS of the successful trials:
%
%   experiments/simulation/              (feeds fig:variable_topo_metrics and the prose bound)
%       familia_c2_pendiente        n=15  mean 1.469  max 1.657
%       familia_c1_terreno_rugoso   n=15  mean 1.586  max 3.070
%   experiments/R3-04-noise-robustness/data/   (feeds tab:fsm_comparison's terrain rows)
%       familia_c2_pendiente        n=15  mean 2.997  max 16.876
%       familia_c1_terreno_rugoso   n=15  mean 1.784  max 1.927
%
% So the prose bound is exactly right FOR ITS OWN dataset: 1.586 is the rugged mean (the "1.59") and
% 1.469 the slope mean (the "1.47" in FigC's caption). The table takes only the three noise-level-0
% trials of the OTHER battery: 1.613, 16.876, 1.764 -> mean 6.751, population SD 7.160 (matches the
% table exactly, confirming the SD is the population one).
%
% Two separate causes, and the second is the structural one:
%
% 1) The 16.876 value is a logging artifact, not an attitude measurement. In that trial
%    (familia_c2_pendiente/test_006_SUCCESS) the accumulated Roll RMS stays below 1.486 deg for the
%    first 67.8 s -- its median over the whole trial, 1.398, matches its two peers (1.423, 1.388) --
%    then climbs 1.466 -> 3.257 -> 9.453 -> 14.275 -> 16.876 in 0.6 s, after which EVERY column of
%    unified_metrics.csv (roll_rms, pitch_rms, Tresponse, Tswitch, latency, lambda) repeats
%    bit-for-bit for 15 samples / 3.0 s. An accumulated RMS cannot stay bit-identical while new data
%    arrives, so the series stopped updating. Its two peers end on their own maximum with no
%    repetition at all (1 identical final sample). The trial also "lasts" 78.6 s vs 75.6 and 76.6 for
%    the peers -- exactly the frozen tail -- and its stability_log.csv is empty (header only).
%    NOTE: familia_c2_pendiente/test_015 (noise level 4, roll_rms 8.136) looks like the same artifact.
%    It reaches no published number, since the table uses only noise level 0, but anyone recomputing
%    this family should check it.
%
% 2) Even with that trial gone the contradiction would survive: the table's Rugged Terrain / Neural
%    row is 1.827 deg, already above the 1.59 bound, with no outlier involved -- simply because the
%    R3-04 battery runs higher on that terrain (mean 1.784 vs 1.586). Fixing only the outlier would
%    not have been enough; the datasets had to be told apart in the text.
%
% WHAT CHANGED (both inside existing \revblue{} this round)
% a) The prose bound is now scoped to the datasets that actually support it, and points the reader at
%    the table instead of leaving the mismatch for them to find:
%    before: ...; the mean Roll RMS stays at or below 1.59 deg and the mean Pitch RMS at or below
%            2.29 deg across these suites and both variable-topography terrains, indicating that ...
%    after:  ...; within those two datasets -- the four scenario suites and the variable-topography
%            battery of 15 trials per terrain -- the mean Roll RMS stays at or below 1.59 deg and the
%            mean Pitch RMS at or below 2.29 deg, indicating that ... locomotion phase. The two
%            terrain rows of Table~\ref{tab:fsm_comparison} come from an independent re-run of those
%            scenarios rather than from this battery, and are not bounded by these values; see that
%            table's caption.
% b) tab:fsm_comparison's caption gained, after the existing N sentence: the terrain rows are an
%    independent re-run (not the Section~\ref{sssec:variable_topo_quant} battery), so their Roll RMS
%    is not comparable with the bounds reported there; and the Inclined Slope/Neural mean is dominated
%    by one of its three trials, whose accumulated Roll RMS stays below 1.49 deg for the first 68 s
%    and then rises abruptly over the final second before the telemetry stops updating, that trial's
%    median (1.40 deg) matching the other two -- which is why the SD here exceeds the mean.
%
% The caption describes the observation (value rises, then stops updating), not a diagnosis of the
% cause, which is inferred and not independently confirmed in Gazebo.
%
% NO NUMBER WAS CHANGED AND NO TRIAL WAS EXCLUDED. 6.751 +/- 7.160 still stands in the table, now with
% the context a reader needs. The alternative considered and rejected by the author was excluding
% test_006 the way familia_a_obstaculo/test_007_SUCCESS (roll_rms 37.32) already is in
% build_master_table.py::EXCLUDED_TRIALS: it would drop that row to N=2 (1.689 +/- 0.076) and, unlike
% the Obstacle case, the spike here sits at the end of an otherwise normal trial, so suppressing it
% would hide a real event rather than a measurement error.
%
% STILL OPEN (not this row): whether the two batteries should be reconciled into one for the final
% version. That touches several tables and figures, and the table's restriction to noise level 0
% exists so the Neural side is comparable with the FSM, which ignores noise entirely.
